\documentclass[usenames,dvipsnames]{beamer}

\mode<presentation> {
\usetheme{Montpellier}
}

\usepackage{graphicx} % Allows including images
\usepackage{booktabs} % Allows the use of \toprule, \midrule and \bottomrule in tables
\usepackage[frenchb]{babel}
\usepackage[T1]{fontenc}
\usepackage[utf8]{inputenc}
\setcounter{tocdepth}{2}
\graphicspath{{../Figures/}}

\DeclareMathOperator*{\argmin}{arg\,min}
\DeclareMathOperator*{\argmax}{arg\,max}
\newtheorem*{proposition}{Proposition}

% exercices comptés
\newcounter{numexos}%Création d'un compteur qui s'appelle numexos
\setcounter{numexos}{0}%initialisation du compteur
\newcommand{\exercice}[1]{%Création d'une macro ayant un paramètre
\addtocounter{numexos}{1}%chaque fois que cette macro est appelée, elle ajoute 1 au compteur numexos
\textcolor{DarkOrchid}{Exercice\,\thenumexos\,:}\,#1%Met en rouge Exercice et la valeur du compteur appelée par \thenumeexos
}

\setbeamertemplate{title page}{
    \centering

    \huge\bfseries
    \large{Solution to the linear regression in dimension 1}

\begin{figure}[!h]
\includegraphics[width=8cm]{optimal_regression}
\end{figure}

    }

\title[FTML]{title}


\date{} % Date, can be changed to a custom date

\makeatletter
\let\@@magyar@captionfix\relax
\makeatother

\begin{document}

\begin{frame}
\titlepage % Print the title page as the first slide
\end{frame}

\begin{frame}{Linear regression}

    Formalization:
    \begin{itemize}
        \item input space (temperature): $\mathcal{X} = \mathbb{R}$
        \item output space (power consumption): $ \mathcal{Y} = \mathbb{R}$
	\item dataset: $D_n=\{(x_1, y_1), \dots, (x_n, y_n), i\in[1, n]\}$.
    \end{itemize}

    When doing linear regression, our estimator is of the form:

    \begin{equation}
        h(x) = \theta x + b
    \end{equation}

    with $\theta\in \mathbb{R}$, $b\in\mathbb{R}$.

\end{frame}

\begin{frame}{Empirical risk minimization}

    \begin{equation}
        R_n(\theta, b) = \frac{1}{n}\sum_{i=1}^{n} (\theta x_i+b-y_i)^2
    \end{equation}

    We want to find $\theta$ and $b$ such that $R_n(\theta, b)$ has the
    \textbf{smallest possible value}.

    Since $(\theta, b) \mapsto R_n(\theta, b)$ is convex, finding its global minimum is equivalent to finding the points where the gradient cancels, which means that the partial derivatives with respect to both $\theta$ and $b$ are equal to $0$.

\end{frame}

\begin{frame}{Derivatives}

    We note that if we look for the $(\theta, b)$ values that cancel the gradient, we can remove the $\frac{1}{n}$ factor without changing the result.
    \begin{equation}
    \begin{aligned}
        \label{eq:}
        \frac{\partial R_n}{\partial \theta}(\theta, b) &= \sum_{i=1}^{n} 2(\theta x_i+b-y_i)x_i\\
        &= 2\Large[\theta \sum_{i=1}^{n} x_i^2 + b\sum_{i=1}^{n} x_i - \sum_{i=1}^{n} x_i y_i \Large]\\
    \end{aligned}
    \end{equation}


    \begin{equation}
    \begin{aligned}
        \label{eq:}
        \frac{\partial R_n}{\partial b}(\theta, b) &= \sum_{i=1}^{n} 2(\theta x_i+b-y_i)\\
        &= 2\Large[\theta \sum_{i=1}^{n} x_i + nb- \sum_{i=1}^{n} y_i \Large]\\
    \end{aligned}
    \end{equation}
\end{frame}

\begin{frame}{}
    Hence we have a system of $2$ equations with $2$ unknowns (dropping the $\theta^*$ notation)

    \begin{equation}
        \theta \sum_{i=1}^{n} x_i^2 + b\sum_{i=1}^{n} x_i - \sum_{i=1}^{n} x_i y_i = 0
    \end{equation}

    \begin{equation}
        \theta \sum_{i=1}^{n} x_i + nb- \sum_{i=1}^{n} y_i = 0
    \end{equation}
\end{frame}

\begin{frame}{}
    Which means

    \begin{equation}
        b =  \frac{1}{n}(\sum_{i=1}^{n} y_i -\theta \sum_{i=1}^{n} x_i)
    \end{equation}

    \begin{equation}
        \theta \sum_{i=1}^{n} x_i^2 + \frac{1}{n}\large(\sum_{i=1}^{n} y_i -\theta \sum_{i=1}^{n} x_i \large)\sum_{i=1}^{n} x_i - \sum_{i=1}^{n} x_i y_i = 0
    \end{equation}

\end{frame}

\begin{frame}
    Finally:

    \begin{equation}
        \theta (\sum_{i=1}^{n} x_i^2 - \frac{1}{n}(\sum_{i=1}^{n} x_i)^2) +\frac{1}{n}\sum_{i=1}^{n} x_i\sum_{i=1}^{n} y_i - \sum_{i=1}^{n} x_i y_i = 0
    \end{equation}

    or

    \begin{equation}
        \theta^* = \frac{\sum_{i=1}^{n} x_i y_i - \frac{1}{n}\sum_{i=1}^{n} x_i\sum_{i=1}^{n} y_i}{\sum_{i=1}^{n} x_i^2 - \frac{1}{n}[\sum_{i=1}^{n} x_i]^2}
    \end{equation}
\end{frame}


\end{document} 
